\section{Discusión}
La reducción a un permanente ocurre después de elegir las últimas filas. El orden uniforme
permite probar la cota geométrica incluso para una solución fija. Bajo la ley adaptativa,
el mismo fenómeno se obtiene a partir de las probabilidades de que sobrevivan filas y columnas
prescritas. Esta distinción explica por qué las dos cotas para el conteo admiten un mismo
último tablero aunque sus historias tengan distribuciones diferentes.

La estimación del permanente es uniforme en una clase con cotas de difusión y de brecha fijas. Su resto
tiende a cero al crecer el tamaño del último tablero. El orden de límites de la
Sección~\ref{balance} permite utilizar precisamente esa información. Para obtener la tasa
$O(1/n)$ de la Conjetura~10.1 de~\cite{pred} sería necesario controlar cuantitativamente
el error terminal y su dependencia del suelo $K$.
